\chapter{What it means for a family to be uniform}\label{tx:ch:uniformity}\label{tx:sec:systems}
The number of variables in our examples grows, but the rule describing
them does not. At universe size $n$ there may be a variable $x_{ij}$ for
every pair of labels $i,j<n$, an equation for every row, and a target
for every distinguished label. The classification theorem requires a
precise version of this fixed-rule condition. Symmetry at each size,
by itself, does not imply the four complexity profiles: the equations
could be chosen differently at each size, according to an arbitrary
rule. We need symmetry together with one coefficient rule which remains
the same as the universe grows.

For the main proof, this chapter has three jobs: specify the labelled
variables, fix the coefficient rule for the axioms, and do the same for
the requested targets. The intervening orbit-sum and compression results
explain the geometry of supported expressions, but the linear lower
bound does not use compression. On a first reading, one may pass from
Definition~\ref{tx:def:template} and its examples directly to
Sections~\ref{tx:sec:target-request}--\ref{tx:sec:target-budgets}, then
return to the structural results later.

\section{Labels, support, and relabelling}
We also call the universe size $n$ the \emph{rank} of the labelled
family. It counts available labels. The dimension of a polynomial
space at that rank is a separate number, usually much larger.

Write $[n]=\{0,\ldots,n-1\}$. Think of its members as names rather
than numerical values. A variable $x_{ij}$ is indexed by the ordered
pair $(i,j)$, and relabelling by a permutation $\pi$ sends it to
$x_{\pi(i),\pi(j)}$. The labels needed to specify that index are
$\{i,j\}$; this is the support of the pair, and therefore of the
variable indexed by it. When $i=j$, its support has size one, not two.
Permutations fixing those labels leave that variable unchanged. We
start with individual indices because monomials and whole polynomial
expressions can have different kinds of support, as the row-sum example
below will show.

\begin{definition}[Bounded basis system]\label{tx:def:basis}
A bounded basis system is a collection of objects labelled from one
infinite pool $\NN$, with an action of label permutations. There is
one fixed integer $b$, independent of $n$ and of the object, and each
object $v$ is assigned a set $\supp(v)\subseteq\NN$ with
$|\supp(v)|\le b$. Supports are equivariant:
$\supp(\pi v)=\pi(\supp(v))$. Every permutation fixing
$\supp(v)$ pointwise fixes $v$. At rank $n$, we keep precisely the
objects whose supports lie in $[n]$; we require that this set be finite.
We call $b$ a support bound for the system.
\end{definition}

Examples include ordered pairs, tuples with specified equality patterns,
and finitely many tagged sorts of such tuples. There may also be objects
with empty support: a distinguished variable $y$ fixed by every label
permutation is an example. A support bound is a bound on the labels
needed by one object, not on the number of objects present at rank $n$.
The important requirement is the \emph{same bound at every rank}.
Finiteness at one rank would not provide that: a bound of $n$ would
grow with the universe and would not be a bounded-support system in
this sense.

A variable basis $\mathcal V$ gives the ring
$K[x_v:v\in\mathcal V_n]$. We write $P_{n,\le D}$ for its vector
space of polynomials of total degree at most $D$. If each variable has support at most $k$,
a monomial of degree at most $d$ involves at most $kd$ labels: it has
at most $d$ variable factors, each involving at most $k$ labels. Powers
and overlapping supports may make the true number smaller.

\begin{example}[The row sum mentions every label]\label{tx:eg:row-support}
With variables $x_{ij}$, consider
\[
                         f_{i,n}=\sum_{j<n}x_{ij}-1.
\]
Its index is $i$, of support one. Every permutation of $[n]$ fixing $i$ fixes the
whole polynomial, because it merely permutes the summands. Nevertheless
the union of labels mentioned in the polynomial is generally all of
$[n]$. Each monomial uses at most two labels, and the index uses one;
neither assertion says that the whole polynomial uses only one or two
labels. This distinction permits totality sums in the theory.
\end{example}

\section{Stable coefficient rules}
\begin{definition}[Uniform axiom template]\label{tx:def:template}
A uniform template consists of a variable basis $\mathcal V$ with
support bound $k\ge1$, an axiom-index basis $\mathcal A$ with support
bound $s$, and coefficients $c(a,\mu)\in K$. Here $a$ is an axiom
index and $\mu$ is a monomial in the labelled variables. The rule is
independent of rank and invariant under simultaneous relabelling:
\[
                 c(\pi a,\pi\mu)=c(a,\mu).
\]
For a fixed $D$, it vanishes whenever $\deg\mu>D$. At rank $n$ the
inputs are
\begin{equation}\label{tx:eq:template-instance}
 \Gamma_n=\{f_{a,n}:a\in\mathcal A_n\},\qquad
 f_{a,n}=\sum_{\mu\text{ supported in }[n]}c(a,\mu)\mu.
\end{equation}
Fix $n_0\ge\max\{1,kD+s\}$ as part of the template data.
\end{definition}

Because there are finitely many variables at each rank and the degree
is bounded, each displayed sum is finite. The same coefficients are
used when the universe grows; there are simply more monomials available.
For the row sum, the coefficient of $x_{ij}$ is one when its row index
is $i$, the constant coefficient is minus one, and other coefficients
are zero. None of these rules asks how large $n$ is.

For a group acting on indices, the \emph{orbit} of $a$ is the set
of all its relabellings $\pi a$. Its \emph{stabiliser} consists of the
permutations which fix $a$. These familiar group-action notions have
a specific use here: indices in the same orbit describe the same
coefficient pattern with different label names. Fixing a support
\emph{pointwise} means fixing each of its labels, not merely preserving
the set of labels as a whole. Later, $H_S$ denotes the subgroup fixing
the distinguished set $S$ pointwise.

The stable range allows us to recognise the degree of an axiom from
its index orbit. A monomial witnessing the largest nonzero coefficient
uses at most $kD$ labels. Together with the axiom index it uses at most
$kD+s$. A permutation fixing the index can move all the additional
labels into $[n]$ when $n\ge kD+s$. Hence, beyond $n_0$, each index
orbit has one fixed polynomial degree. This matters when deciding which
multiplier--axiom products are legal at a proposed proof degree.

\begin{example}[Symmetric but not uniform]\label{tx:eg:not-uniform}
Choose an arbitrary set $A$ of positive integers. At rank $n\in A$
let the inputs be $\{1\}$, and at other ranks let there be no inputs.
Each instance is invariant under every label permutation. Yet the
refutation degree is zero on $A$ and infinite elsewhere. Arbitrary
choices of $A$ cannot satisfy the periodic conclusions of this book.
They are excluded because the axiom rule changes with rank.
\end{example}

\section{Supported expressions are orbit-sum expressions}\label{tx:sec:supported-expressions}
The row sum suggests a converse question. Suppose we are given a
polynomial and know that every permutation outside a short list of
labels fixes it. Can we describe the polynomial by a uniform rule?
At one rank the answer is yes. To represent a whole family by one
rule, we also need compatibility as rank changes.

Here is the support notion for a whole expression. At rank $n$, a set
$S\subseteq[n]$ \emph{supports a polynomial $F$} if every permutation
of $[n]$ fixing $S$ pointwise fixes $F$. The permutation may move every
other label. This is different from collecting every label mentioned
in its monomials. For example, $\sum_{j<n}x_j$ has empty support as an
expression, although its monomials collectively mention all labels.
The row sum $f_{i,n}$ is supported by $\{i\}$ in this sense.

\begin{proposition}[Orbit-sum description]\label{tx:prop:orbit-sum}
Let $F$ have degree at most $D$, and let $H_S$ be the permutations of
$[n]$ fixing $S$ pointwise. The polynomial $F$ is supported by $S$
exactly when it has a unique expression
\[
             F=\sum_O a_O\,\Omega_{O,n},\qquad
             \Omega_{O,n}=\sum_{\mu\in O}\mu,
\]
where $O$ runs over the $H_S$-orbits of degree-at-most-$D$ monomials.
\end{proposition}
\begin{proof}
Relabelling permutes the monomial basis. Thus $hF=F$ for all $h\in H_S$
exactly when monomials in the same orbit have the same coefficient.
Collect these equal coefficients to obtain the displayed sum.
Different orbits have disjoint monomial sets, so the orbit sums are
linearly independent and their coefficients are unique.
\end{proof}

This is a grouping of distinct monomials, not an average over a group.
In particular, it is valid in characteristic $p$ even when the order
of the stabiliser is divisible by $p$.

\begin{example}[All supported linear expressions in one-variable indices]\label{tx:eg:supported-linear}
For variables $x_j$ and a distinguished label $i$, every linear
polynomial fixed by the permutations fixing $i$ has the form
\[
                    a x_i+b\sum_{j\ne i}x_j+c.
\]
For $n\ge2$ there are three monomial orbits: the constant monomial, $x_i$, and all
other variables. At $n=1$ the last orbit is absent, so its coefficient
cannot yet be recovered from the polynomial. At $n\ge2$ all three
orbits are present. Keeping $a,b,c$ fixed gives a uniform expression
at every larger rank.
\end{example}

\subsection{From one supported expression to one uniform template}
Return to a variable basis of support at most $k$. Each monomial of
degree at most $D$ uses at most $r=kD$ labels. Let $s=|S|$ and choose
an ordering $a_0$ of $S$. Relative to those $s$ distinguished labels,
every global monomial orbit can be represented using at most $s+r$
labels in total. Hence all its orbit types are visible at any rank
$n\ge s+r$ containing $S$.

\begin{proposition}[Uniform extension of one supported polynomial]\label{tx:prop:supported-extension}
Suppose $n\ge s+kD$ and $F\in P_{n,\le D}$ is supported by an
$s$-element set $S$. Using ordered distinct $s$-tuples as indices, there
is a unique rank-independent, simultaneously relabelling-invariant
coefficient rule of degree at most $D$ whose polynomial at index
$a_0$ and rank $n$ is $F$. It gives a uniform axiom or target orbit
with index-support bound $s$ and no increase in polynomial degree.
\end{proposition}
\begin{proof}
Consider a pair $(a,\mu)$ consisting of an ordered distinct $s$-tuple
and a monomial of degree at most $D$. First relabel $a$ to $a_0$.
The remaining monomial labels can then be moved into $[n]$ by a
permutation fixing $S$, because there are at most $kD$ of them and
$n-s\ge kD$. Thus every global orbit of pairs has a representative
$(a_0,\mu_0)$ at the chosen rank.

Assign to that orbit the coefficient of $\mu_0$ in $F$. This value
does not depend on the representative. Indeed, a global relabelling
between two such representatives fixes $a_0$. Its restriction to the
labels used by the tuple and monomial is a finite injection within
$[n]$, which extends to a permutation of $[n]$ fixing $S$. It therefore
relates the two monomials by $H_S$, and their coefficients agree by
Proposition~\ref{tx:prop:orbit-sum}.

The assigned value defines $c(a,\mu)$ on all pairs, with zero
coefficients above degree $D$. Being constant on simultaneous pair
orbits makes it relabelling invariant. Summing its coefficients over
rank-$m$ monomials gives the required polynomial for every index
available at rank $m$. At the original rank and index, every coefficient
agrees with $F$. All pair orbits already occurred there, so the
coefficient rule is unique.
\end{proof}

Thus one supported polynomial at a sufficiently large rank determines
its uniform continuation. This does not say that a separately chosen
polynomial at every rank necessarily follows that continuation.

\subsection{The additional condition for a whole family}
For $m\le n$, let $\rho_{n,m}$ retain the coefficients of monomials
supported in $[m]$ and discard the other monomials. This is the
coefficient restriction map, rather than a transport rule for proof
steps. A stable uniform family satisfies
\begin{equation}\label{tx:eq:coefficient-coherence}
        \rho_{n,m}(F_{a,n})=F_{a,m}
        \quad\text{whenever the index }a\text{ is supported in }[m].
\end{equation}

\begin{proposition}[Characterisation of a stable coefficient rule]\label{tx:prop:coherent-family}
Fix bounded variable and index bases and a degree bound $D$. A family
$F_{a,n}\in P_{n,\le D}$, defined for every $n\ge1$ and every index
$a$ available at that rank, comes from one uniform coefficient rule if
and only if it is equivariant under permutations at each rank and
satisfies \eqref{tx:eq:coefficient-coherence}.
\end{proposition}
\begin{proof}
A uniform rule gives equivariance by simultaneous relabelling of
indices and monomials. Restriction preserves exactly its coefficients
on the retained monomials, proving coherence.

Conversely, for a global pair $(a,\mu)$ choose a rank containing all
its labels and define its coefficient to be the coefficient of $\mu$
in $F_{a,n}$. Coherence makes this independent of the chosen rank.
To compare it with the coefficient of $(\pi a,\pi\mu)$, choose a
larger rank containing the labels of both pairs. The action of $\pi$
on the finite support of the pair extends to a permutation of that
rank. Equivariance there gives equality of the two coefficients.
The resulting rule has all the required properties, and reproduces
the family coefficient by coefficient.
\end{proof}

For example, $F_n=a_n\sum_{i<n}x_i$ is invariant at every rank for
\emph{any} sequence of scalars $a_n$. If the scalars change, the
coefficient of a retained variable changes, and coherence fails.
Small support and bounded degree control the form of an expression;
coherence controls its behaviour across different universe sizes.

\section{Smaller-support generators: what can be compressed?}\label{tx:sec:support-compression}
The bound $kD$ applies to one monomial. It need not bound the number
of distinguished labels needed to specify a polynomial. A further,
different possibility is to replace a polynomial's entire permutation
orbit by generators with smaller support. First we state a sufficient
condition for replacing axioms without changing proof degrees.

\begin{proposition}[Degree-preserving axiom replacement]\label{tx:prop:filtered-replacement}
Suppose two axiom sets $\Gamma$ and $\Gamma'$ in the same polynomial
ring satisfy, for every integer $e\ge0$,
\[
 \Span_K\{f\in\Gamma:\deg f\le e\}
       =\Span_K\{g\in\Gamma':\deg g\le e\}.
\]
Then every target has the same least NS degree and the same least PC
degree from the two sets, including degree zero and infinity.
\end{proposition}
\begin{proof}
An input $f\in\Gamma$ of degree $e$ is a constant-coefficient linear
combination of inputs from $\Gamma'$ of degree at most $e$. In an NS
certificate replace $f$ by that combination. Every new multiplier--input
product has degree no larger than the product it replaces, since its
new input degree is at most $e$. This transforms any certificate
without increasing its degree. In a PC derivation, derive each old
input by these linear combinations of new inputs, within its old
degree, and then reuse the original inference steps. This also
preserves the degree budget. Interchanging the two sets gives the
converse transformations and proves equality of their least degrees.
\end{proof}

The increasing spaces
$A_e=\Span_K\{f\in\Gamma:\deg f\le e\}$ form a \emph{filtration}:
$A_0\subseteq A_1\subseteq\cdots$, organised by the usable input degree.
The replacement criterion compares this entire filtration.
The unrestricted span forgets which inputs were usable at each cutoff.

An arbitrary change of basis for the full axiom span need not meet this
condition. In Example~\ref{tx:eg:cancellation}, replacing $f_2$ by
$f_1+f_2=x+1$ leaves the axiom span unchanged, but now
$z=z(x+1)+f_3$ is an NS certificate of degree two. From the original
axioms the least NS degree is three, as Exercise~\ref{tx:ex:small-gap}
shows. Equality of the unrestricted span alone therefore does not
justify a claim about exact proof degrees. The proposition supplies
a sufficient condition by comparing the inputs usable at \emph{each}
degree.

The next example satisfies the criterion and supplies a successful
support compression.

\begin{proposition}[Compression of a linear subset sum]\label{tx:prop:linear-compression}
Let $\varnothing\ne S\subsetneq[n]$, put $s=|S|$, and consider
$F_S=\sum_{i\in S}x_i$. Its smallest support as an expression has size
$\min\{s,n-s\}$. Nevertheless its $S_n$-orbit span has generators
supported on at most two labels. More precisely, with $s$ interpreted
as a scalar in $K$,
\[
 \Span_K\{\pi F_S:\pi\in S_n\}=
 \begin{cases}
     \Span_K\{x_i:i<n\},&s\ne0\text{ in }K,\\
     \Span_K\{x_i-x_j:i,j<n\},&s=0\text{ in }K.
 \end{cases}
\]
\end{proposition}
\begin{proof}
The coefficient of $x_i$ is one on $S$ and zero off $S$. If labels
outside a proposed support include one label of each kind, swapping
those two labels changes the polynomial. Thus a support must contain
all of $S$ or all of its complement, and either choice is sufficient.
This proves the minimum-support formula.

Take $i\in S$ and $j\notin S$. Then
\[
                       F_S-(ij)F_S=x_i-x_j.
\]
Relabelling shows that the orbit span contains every difference.
Those differences span the subspace
$A=\{\sum_i a_ix_i:\sum_i a_i=0\}$: fix a label $j$ and write a
vector of $A$ as $\sum_{i\ne j}a_i(x_i-x_j)$.
Each orbit vector has coefficient sum $s$. If $s=0$ in $K$, the
orbit span is contained in $A$ and is therefore exactly $A$. If
$s\ne0$, it also contains $F_S\notin A$; since $A$ has codimension
one, it is the whole linear polynomial space.
\end{proof}

For example, $\sum_{i\in S}x_i$ can need many labels to specify even
though each monomial uses only one. The proposition replaces its orbit
span, not the polynomial itself, by one- or two-label generators.
All replacements are linear combinations of the original orbit
polynomials and conversely. Thus making this replacement in an axiom
set changes neither the NS nor the PC consequence spaces at any
positive degree budget: the starting linear spans agree, and all
replacement equations have degree one.

The two-label bound can really be necessary. Suppose the characteristic
$p$ divides both $n$ and $s$, with $n\ge3$. A vector of $A$ supported
by at most one label $i$ has the form
$a x_i+b\sum_{j\ne i}x_j$. Its coefficient sum is
$a+(n-1)b=a-b$, so $a=b$. Such vectors are scalar multiples of the
fully invariant sum $\sum_jx_j$. They span only one dimension, whereas
$A$ has dimension $n-1$. Consequently one-label generators cannot
replace the whole orbit span in this case.

\begin{example}[Four-label axioms replaced by two-label axioms]\label{tx:eg:four-to-two}
In characteristic two, let $\Gamma_n$ consist of all sums
$\sum_{i\in S}x_i$ over four-element subsets $S\subseteq[n]$.
For every $n\ge5$, their span is exactly the span of the differences
$x_i-x_j$. Thus we can replace all the four-label equations by the
two-label equations $x_i-x_j=0$, preserving every NS and PC consequence
degree. Both descriptions have fixed coefficient rules across ranks.
The starting-rank condition matters: at $n=4$ there is only one
four-element subset, so its single sum does not span all differences.

At $n=12$, an individual four-label sum has smallest support four,
whereas its orbit span has two-label generators. Since the characteristic
divides both 12 and 4, the preceding argument shows that one-label
generators would not suffice. This gives a concrete reason to distinguish
the polynomial's own support from the support of generators for its
orbit span.
\end{example}

The general compression question has been resolved in separate
companion work~\cite{CompressionCompanion}; its proof is not developed
in this book. In the stable-template setting over a field, put $r=kD$.
There is one coherent replacement template with the same variables
and degree bound, and with axiom-index support at most $2r$, such that
for every
\[
                  n\ge\max\{n_0,\max(s,r)+2r\}
\]
the old and new axiom spans agree at \emph{every input-degree cutoff}.
The replacement criterion above then preserves every target's least
NS degree and its least ordinary-PC degree separately, including zero
and infinity. This is an optional normal form, not an assertion about
arbitrary submodules without a bound on the starting supports, and not
an executable compression algorithm. No equivalence at the exceptional
small ranks is claimed.

The monomial bound $kD$ alone does not establish that compression;
the companion proof supplies an additional structural argument.
The lifting and classification proofs below retain the given index
support $s$ in their budget $2kd+s+t$; they do not assume that it has
already been compressed to $kD$ or $2kD$.

\section{Targets also need a fixed rule}\label{tx:sec:target-request}
\begin{definition}[Uniform target request]\label{tx:def:target}
A uniform target family for a fixed axiom template has a target-index
basis $\mathcal B$ with support bound $t$, a fixed degree bound $D_G$,
and coefficients $\gamma(b,\mu)\in K$. They are independent of rank,
invariant under simultaneous relabelling of $b$ and $\mu$, and zero
above degree $D_G$. The rank-$n$ target indexed by $b\in\mathcal B_n$ is
\[
                  G_{b,n}=\sum_{\mu\text{ supported in }[n]}
                                     \gamma(b,\mu)\mu.
\]
For $S=\NS$ or $\PC$, define the all-target profile
\begin{equation}\label{tx:eq:target-profile}
 h_S(n)=\max_{b\in\mathcal B_n}h_S(\Gamma_n\vdash G_{b,n}).
\end{equation}
The empty maximum is zero. If any requested target is underivable,
the maximum is infinity.
\end{definition}

The statement $h_S(n)\le d$ means that \emph{every} requested target
has some derivation within degree $d$. Those derivations need not share
lines, and their total length is not bounded. This is a maximum of
individual minimum degrees, not the degree of the sum of the targets.
Nor are the targets inserted into the axiom system.

There are three scopes which we will keep distinct.
\begin{enumerate}
\item A fully invariant target has a single index of empty support.
Its polynomial $G_n$ is fixed by every permutation of $[n]$. For
example, $G_n=\sum_{i<n}x_i$ is invariant and has degree one, although
it mentions all labels.
\item A fixed target orbit can contain many indices at a rank. For
$G_{i,n}=x_i$, the indices $i$ form one orbit. Any one representative
has the same least degree as every other representative, since
relabelling transports proofs.
\item An all-target request may contain several orbits. Its profile
is the maximum of their individual profiles.
\end{enumerate}

Every global index orbit has a representative at rank $n\ge t$.
Indeed, move its support into $[t]$. There are only finitely many global
orbits because $\mathcal B_t$ is finite. Within one orbit, two rank-$n$
indices can be related by a permutation of $[n]$: first use the given
relabelling on their supports and then extend that finite bijection.
Thus choosing a representative which varies with $n$ is harmless
\emph{provided it stays in one fixed orbit}. Switching between different
orbits according to an arbitrary rule is not a uniform request.

\section{Stable target degrees and the label budgets}\label{tx:sec:target-budgets}
The same support argument used for axioms applies to targets, now with
at most $kD_G+t$ labels. We therefore use
\begin{equation}\label{tx:eq:stable-rank}
                          n_*=\max\{n_0,1,kD_G+t\}.
\end{equation}
Beyond this rank, every target orbit has a fixed degree. Their largest
degree is a floor for both all-target profiles. A proof budget below
that floor fails; we do not replace the target by its lower-degree
part in order to make it fit.

For an axiom multiplier, the combined support of a monomial of degree
at most $d$ and an axiom index is at most $kd+s$. For multiplication
inside a PC proof, a variable and a monomial of degree at most $d-1$
use at most $k+k(d-1)=kd$ labels. These elementary counts explain why
proof degree enters the transfer theorem linearly.

General indexed targets require a further detection argument. The
resulting safe budget is
\begin{equation}\label{tx:eq:support-budget}
                R(d)=2kd+s+t=\kappa d+\sigma,
                    \qquad\kappa=2k,\quad\sigma=s+t.
\end{equation}
For one fully invariant target the sharper budget is $R(d)=kd+s$.
The factor of two in the general case is not an extra degree cost in a
proof. It is a support cost in the finite-dimensional argument which
shows that target membership is reflected, as well as preserved, when
labels are deleted.

\section{Where first-order systems enter}
A universal clause is a disjunction of literals required to hold for
every assignment of its logical variables to universe labels. A
\emph{literal} is an atom or its negation. A \emph{grounding} replaces
all those logical variables by particular labels; one clause pattern
therefore produces many polynomial equations at rank $n$.
For example, grounding $R(u)\vee\neg R(v)$ at $u=i,v=j$ gives the
clause whose polynomial translation appears below.

The normalised syntax is fixed before proof degrees are measured.
\emph{Flattened} function-graph atoms name an output explicitly,
as in $f(i)=j$, instead of putting nested function terms inside an
atom. They can therefore be indexed by bounded tuples of labels.
Equality between two grounded labels is evaluated as equality of
those labels; it is not an extra unknown relation.

A fixed first-order sentence can supply a uniform template. Relation
atoms receive Boolean variables. A function $f$ is represented by its
graph variables $x^f_{\bar i,j}$, with equations
\[
 \sum_{j<n}x^f_{\bar i,j}-1=0,\qquad
 x^f_{\bar i,j}x^f_{\bar i,j'}=0\quad(j\ne j').
\]
Boolean equations make each value zero or one; collision equations give
at most one output; the sum gives at least one. The collisions cannot
be dropped in characteristic $p$, where $p+1$ ones also sum to one.
A clause is translated as a product of its literal-falsity indicators.
For example $R(i)\vee\neg R(j)$ gives $(1-r_i)r_j=0$.

Fixed arities and a fixed finite list of clause patterns give bounded
variable and axiom-index supports, and expansion gives a stable
coefficient rule. Functions of arity zero represent constants. There
is no built-in numerical order on the labels: an order, if used in a
sentence, is a relation symbol with its own axioms. Later examples keep
this graph-and-sum encoding visible, since an arbitrary change of
encoding need not preserve minimum degree.

\begin{proposition}[Normalised first-order compiler]\label{tx:prop:compiler}
A fixed finite list of universal clauses in flattened relation, function
graph, and equality atoms gives a uniform template by the displayed
translation. At every $n\ge1$, its roots over $K$ correspond exactly
to models on $[n]$ of that normalised sentence.
\end{proposition}
\begin{proof}
Use atom tuples as variable indices, and grounding tuples as clause
indices. Totality and collision indices are fixed tuples as well.
Expansion of a clause product depends only on the fixed syntax and on
equalities among the grounding labels; it therefore gives stable,
simultaneously invariant coefficients, including repetitions and
cancellations. If $a_R,a_f$ are the maximum relation and function
arities, safe parameters are
\[
 k=a_R+a_f+1,\qquad s=k+1+\max_C\operatorname{nvars}(C),\qquad
 D=2+\max_C|C|,\qquad n_0=kD+s+1.
\]
Here $\operatorname{nvars}(C)$ is one more than the largest variable
index occurring in $C$, and is one if none occurs; an unused variable
slot is harmless. Maxima over empty symbol or clause lists may be taken
as zero where appropriate.

At a root, every relation and graph variable is zero or one. Totality
and collisions select exactly one output for each function input, and
a clause product vanishes precisely when one of its literals is true.
Equality atoms use their actual label-equality values. This constructs
a model. Conversely a model supplies the corresponding Boolean values
and satisfies all translated equations.
\end{proof}

Normalisation is part of the input convention. This proposition does
not say that different normalisations, or arbitrary introductions of
witness functions, preserve exact proof degree. The later examples use
specific normalised sentences so that their equations are unambiguous.

\section{Exercises}
\begin{exercise}\label{tx:ex:support-count}
For variables $x_{ij}$, find the label supports and total degrees of
$x_{ij}^3$ and $x_{ij}x_{jk}$ when $i,j,k$ are distinct. Compare the
actual supports with the general bound $2d$. Explain why the row sum
in Example~\ref{tx:eg:row-support} has an index of support one.
\end{exercise}
\noindent{\small\hyperref[tx:ans:support-count]{Model answer (p.~\pageref*{tx:ans:support-count}).}}\par\medskip

\begin{exercise}\label{tx:ex:orbit-switch}
Take no axioms, a single variable $x$ of empty support, and two target
indices of empty support requesting $0$ and $x$. Determine the two
fixed-orbit profiles and the all-target profile. What happens if one
chooses the first target at sizes in an arbitrary set and the second
elsewhere?
\end{exercise}
\noindent{\small\hyperref[tx:ans:orbit-switch]{Model answer (p.~\pageref*{tx:ans:orbit-switch}).}}\par\medskip

\begin{exercise}\label{tx:ex:stable-floor}
Take variables $x_i$, no axioms, and targets $x_i-x_j$ indexed by
ordered distinct pairs. Use $k=1,s=0,t=2,D_G=1,n_0=1$ to compute $n_*$.
At which positive ranks does this request actually contain a target?
Why is the stated stable rank allowed to be larger than necessary?
\end{exercise}
\noindent{\small\hyperref[tx:ans:stable-floor]{Model answer (p.~\pageref*{tx:ans:stable-floor}).}}\par\medskip

